\section{Zastosowanie Large Language Models do automatycznego code review - stan wiedzy}
\label{sec:llm_acr_sota}

Autorzy współczesnych prac naukowych w dziedzinie informatyki, a konkretnie tematów związanych z przeglądem kodu, wyraźnie idą w stronę
rozwiązań implementujących duże modele językowe (LLM). Podejście to, w odróżnieniu od wcześniejszych podejść
opartych o reguły i klasyczne uczenie maszynowe, pozwala formułować uwagi w języku naturalnym. Warto
zauważyć, że LLM operują na reprezentacjach kodu, co czyni je atrakcyjnymi w przypadkach wymagających zarówno
identyfikacji problemu, jak i uzasadnienia i wskazania lepszego podejścia lub co poprawić.
Stan wiedzy naukowej pokazuje, że skuteczność LLM w automatycznym code review uzależniona jest od dostępnego kontekstu,
sposobu sterowania modelem, integracji z zewnętrznymi danymi i narzędziami analitycznymi oraz metody
ewaluacji jakości wyników.

Literatura wyraźnie podkreśla, że wzrost wyników metryk automatycznych nie wystarcza do oceny dojrzałości
metody Automated Code Review (ACR). Tufano i in. przeprowadzili ręczną inspekcję łącznie 2291 predykcji i zauważyli,
że pomijając analizę jakościową, łatwo pomylić realne postępy z poprawą
w "łatwych" przypadkach. Autorzy stwierdzili także, że zbiory danych zawierać mogą błędy mapowania komentarz-zmiana, co
zniekształca wnioski eksperymentalne \cite{Tufano2024CodeReviewAutomation}. Badania benchmarkingowe pokazują, że
nawet przy obiektywnej weryfikacji poprzez testy jednostkowe skuteczność modeli może być umiarkowana. Cihana i in.
przeprowadzili eksperyment, w którym GPT-4o osiąga 68.50\% trafności klasyfikacji poprawności na zbiorze 492 fragmentów kodu,
a jednocześnie odnotowano niezerowe ryzyko regresji (pogorszenia poprawnego kodu) na poziomie 10.43\% \cite{Cihan2025EvaluatingLLMsForCodeReview}. 
Wyniki te wspierają tezę, że pełna automatyzacja procesu pozostaje ryzykowna i uzasadniają podejścia \emph{human-in-the-loop}.

Jednocześnie rozwijają się prace traktujace code review jako zadanie wielowymiarowe
i wskazujące ograniczenia komentarzy wygenerowanych w sposób prosty przez jeden model. Ren i in. przeanalizowali
30\,157 komentarzy review oraz 29\,612 zmian z 1000 najpopularniejszych repozytoriów na platformie GitHub, wyprowadzili
taksonomię 19 wymiarów w kontekście recenzji. Autorzy dowodzą, że częstym problemem jest niepoprawność sugestii i zbyt duży poziom ogólności
- : 33.2\% sugestii generowanych przez wybrane metody (m.in. 
CodeReviewer i LLaMA-Reviewer) oceniono jako niepoprawne, a zbyt duża ogólnikowość dotyczyła 61.0\% komentarzy 
metod konwencjonalnych wobec 16.4\% w przypadku ChatGPT \cite{Ren2025HydraReviewer}. W ramach tworzenia narzędzia do
automatycznego przeglądu kodu oznacza to, że oprócz samego wygenerowania komentarza, istotne jest zapewnienie sposobu
na jego zweryfikowanie oraz określenie użyteczności i jednoznaczności w procesie MR/PR. Idealnie pasuje to do założeń
niniejszej pracy zawierającymi integracje z platformami Git oraz procesami ciągłej integracji (CI).

Patrząc na to, jak takie rozwiązania wyglądają w praktyce, widać, że kluczem do sukcesu 
jest wpięcie automatycznego review bezpośrednio w codzienne narzędzia zespołu. Nie chodzi tu o tworzenie 
nowych platform, ale raczej o integrację z tym, co już jest używane — czyli 
komentarzami w narzędziach do przeglądu czy samym edytorem, wykorzystując przy tym strukturę kodu. Dobrym 
przykładem jest tutaj podejście opisane w raporcie firmy Ericsson \cite{Ramesh2025EricssonExperience}. Zbudowali oni lekki pipeline 
działający z Gerritem i VS Code, który dodatkowo wspiera się Tree-sitterem (AST), żeby wyniki 
analizy trafiały prosto do typowego flow pracy programisty.

Taką samą logikę można z powodzeniem zastosować na popularnym GitHubie. Automatyczne komentarze mogą 
lądować bezpośrednio w Pull Requestach, co sprawia, że stają się one naturalnym elementem 
przeglądu kodu, a nie tylko zewnętrznym raportem. Od strony technicznej sprowadza się to do 
odpalenia analizy w ramach CI, wygenerowania uwag inline do diffa i — co ważne 
— zostawienia ostatecznego głosu recenzentowi. Takie podejście \emph{human-in-the-loop} jest zgodne z tym, co 
sugeruje literatura \cite{bacchelli2013, Ramesh2025EricssonExperience}, podkreślająca czysto pomocniczy, narzędziowy charakter MCR. Cała ta komunikacja i 
wyciąganie potrzebnych danych opiera się tu po prostu na wykorzystaniu interfejsów API oraz odpowiednio 
skonfigurowanych webhookach.

Kolejne sekcje tego podrozdziału poświęcone są przedstawieniu stanu badań naukowych w pryzmacie konkretnych nurtów technicznych.
Są one odpowiedziami na konkretne ograniczenia:
\begin{itemize}
    \item podejścia oparte na samym promptowaniu \cite{Pornprasit2024FineTuningPromptingCR},
    \item dostrajanie modeli do danych \emph{code review} \cite{Pornprasit2024FineTuningPromptingCR},
    \item metody wykorzystujące zewnętrzny kontekst \cite{Ren2025HydraReviewer},
    \item integracja LLM z analizą statyczną i innymi narzędziami analizy kodu \cite{Jaoua2025LLMStaticAnalyzersCR, Ramesh2025EricssonExperience}.
\end{itemize}

Do sformuowania pełnej luki badawczej dla niniejszej pracy konieczne jest spojrzenie dodatkowo na kwestie problemu
ewaluacji jakości ACR. Poza metrykami leksykalnymi i dopasowania (np. BLEU) w literaturze pojawiają się propozycje
benchmarków nastawionych na zrozumienie zmiany w kontekście jej lokalizacji, np. CodeReviewQA (900 kuratorowanych przykładów z 199 repozytoriów 
i 9 języków). Ograniczają one wpływ "zgadywania" oraz zmniejszają szansę na zanieczyszczenie danych \cite{Lin2025CodeReviewQA}. 
Przegląd literatury w tej postaci pozwoli na krytyczne zidentyfikowanie luki badawczej i potwierdzi zasadność przyjętej stretegii:
wykorzystanie LLM w roli wspierającej recenzenta, z kontrolą jakości poprzez kontekst z repozytorium, komponenty hybrydowe oraz metody ewaluacji bliższe
prawdziwemu procesowi recenzji kodu.

Na podstawie literatury widać, że lukę badawczą można zauważyć w dwóch aspektach - praktycznym i systemowym: 
dominują prace skoncentrowane na pojedynczych elementach (np. generacja komentarzy), 
podczas gdy brakuje rozwiązań \emph{end-to-end}, które działają w realnym workflow PR/MR, 
integrują kontekst repozytorium oraz mają jasno określoną warstwę weryfikacji 
wyników \cite{Tufano2024CodeReviewAutomation, Ren2025HydraReviewer, Ramesh2025EricssonExperience}. 
W ramach niniejszej pracy luka ta będzie wypełniana poprzez zaprojektowanie i opisanie 
systemu ACR, który: (i) działa jako krok CI w PR/MR, (ii) wykorzystuje RAG do 
dostarczania kontekstu projektowego, (iii) łączy generację LLM z warstwą 
kontroli jakości oraz (iv) pozostawia decyzję po stronie recenzenta 
(\emph{human-in-the-loop}) \cite{Meng2025RARe, Jaoua2025LLMStaticAnalyzersCR, Lin2025CodeReviewQA}.

\subsection{Podejścia oparte na promptowaniu (in-context learning)}
\label{subsec:llm_acr_prompting}

Najprostszy sposób na użycie LLM w automatycznym code review to potraktowanie modelu jako mechanizmu
"wnioskowania z kontekstu" (\emph{in-context learning}). W praktyce oznacza to, że wagi modelu nie są zmieniane, a
sterowanie odbywa się poprzez sam prompt: instrukcję, format odpowiedzi i kilka przykładów (albo ich brak - tzw. zero-shot).
Jakość wyników zależy głównie od tego, co i jak trafi do treści promptu - od doboru przykładów/demonstracji (\emph{few-shot}),
ale też od tego, ile kontekstu dołożymy i jak bardzo jest on czysty (np. czy nie ma szumu, sprzecznych wskazówek, niestandardowych 
formatów). Z punktu widzenia wdrożenia takich systemów jest to dobra wiadomość: nie jest potrzebna infrastruktura treningowa, 
możliwe jest szybkie iterowanie. Jednak literatura dość konsekwentnie pokazuje, że takie 
podejście bywa kapryśne i wrażliwe na detale, czasem wręcz na miniaturowe zmiany w promptcie.

Pornprasit i Tantithamthavorna opisują prompting jako alternatywę dla dostrajania, zwłaszcza w obliczu
problemów z danymi na starcie (tzw. \emph{cold-start}). Porównują kilka wariantów, m.in. zero-shot 
i few-shot, a także podejście z personą i bez persony. Ciekawy i praktyczny jest u nich wątek automatycznego doboru
przykładów: few-shot nie jest "ręcznie wybrane", tylko oparte o wyszukiwanie podobnych przypadków metodą BM25.
Jednocześnie liczba przykładów jest ograniczona do trzech, dzięki czemu określić to można jako świadomy kompromis
pomiędzy kosztem tokenów a stabilnością odpowiedzi \cite{Pornprasit2024FineTuningPromptingCR}. Wyniki autorów
zawierają również nieoczywiste wnioski: rozbudowane, wielokrokowe instrukcje nie zawsze pomagają, a w części ustawień lepiej 
działały prompty prostsze, ale o konsekwentnej strukturze \cite{Pornprasit2024FineTuningPromptingCR}. Dla projektu
narzędzia do automatycznego code review jest to dobry argument, aby budować szablony promptów
(np. dla każdego typu plik lub języka programowania) i testować je w kontrolowany sposób, zamiast dołączać co raz następne reguły
bo tak podpowiada logika.

Haider i in. patrzą na prompting w zadaniu generowania komentarzy recenzenckich (RCG) z szerszej perspektywy.
Autorzy zestawiają \emph{few-shot prompting} dla modeli zamkniętych z podejściami opartymi 
o dostrajanie (QLoRA) dla modeli otwartych. W samym wariancie promptingowym też pojawia 
się BM25 do wyboru przykładów, co wzmacnia wniosek, że w praktyce "promptowanie" 
często oznacza ustrukturyzowany ICL (\emph{in-context learning}) wsparty wyszukiwaniem, a nie tylko ręcznie napisany prompt 
\cite{Haider2024PromptingFineTuningRCG}. U Haider i in. ważny jest także kolejny szczegół: relacja między dodatkowymi metadanymi 
a jakością nie jest liniowa i jednoznaczna. Dokładanie np. grafu wywołań funkcji czy streszczenia 
kodu nie zawsze poprawiało rezultaty; efekt zależał między innymi od tego, jak 
dużo kontekstu model realnie potrafi sensownie wykorzystać \cite{Haider2024PromptingFineTuningRCG}.
Wniosek inżynierski na podstawie powyższych obserwacji to: wzbogacanie promptu trzeba traktować jako decyzję projektową i sprawdzać wpływ 
ablacją oraz pomiarem, zamiast przyjmować założenie "im więcej kontekstu, tym lepiej".

Kontekst w promptowaniu ma jeszcze jeden wymiar: nie tylko "kod", ale też 
opis zadania i intencja. Cihan i in. pokazują to na przykładzie procedury 
podobnej do realnego procesu recenzji zmian: model najpierw klasyfikuje poprawność kodu, a 
potem - dla przypadków uznanych za błędne - generuje poprawkę, którą weryfikuje 
się testami jednostkowymi. W ich ustawieniu dorzucenie opisu problemu do promptu poprawiało 
wyniki: dla zbioru 492 bloków kodu GPT-4o osiągnął 68.50\% trafności klasyfikacji poprawności 
oraz 67.83\% skuteczności poprawek, ale dalej występowało ryzyko regresji na poziomie 10.43\% 
\cite{Cihan2025EvaluatingLLMsForCodeReview}. To jest o tyle istotne, że nawet przy sensownie zaprojektowanym promptcie 
nie znika problem "złych zmian". W automatycznym \emph{code review} nie powinno się 
się tego zostawić bez zabezpieczeń: walidacja testami, analiza statyczna czy reguły blokujące 
wciąż wyglądają na konieczny element całego pipeline'u.

W tle tych prac przewija się też temat oceny jakości. Tufano i 
in. zwracają uwagę, że w zadaniach generatywnych (w tym w automatyzacji code 
review) łatwo wpaść w pułapkę metryk czysto ilościowych, a dodatkowo problemy w 
danych (np. błędne mapowania komentarz--zmiana) potrafią sztucznie zawyżać albo zaniżać wyniki \cite{Tufano2024CodeReviewAutomation}. 
Przy promptowaniu ma to szczególne znaczenie, bo iteracje są szybkie i często 
różnice są "na granicy": coś wygląda lepiej w jednej metryce, gorzej w 
innej, a jakościowo bywa różnie. W praktyce to sugeruje łączenie metryk automatycznych 
z oceną jakościową i testami możliwie bliskimi realnemu procesowi recenzji, a nie 
poleganie na jednym wskaźniku.

W efekcie prompting jest sensownym punktem startu dla automatycznego \emph{code review}, zwłaszcza 
gdy danych jest mało i trzeba szybko zbudować prototyp. Jednocześnie z literatury 
da się wyciągnąć kilka dość konkretnych zasad: przykłady warto dobierać automatycznie (np. 
BM25) i kontrolować ich liczbę; kontekst lepiej dobierać selektywnie i 
sprawdzać jego wpływ ablacją; a ocenę wyników oprzeć nie tylko o dopasowanie 
tekstu, ale też o mechanizmy, które łapią błędy merytoryczne i ryzyko regresji 
(testy, statyczna analiza, reguły). To nie rozwiązuje wszystkiego, ale przynajmniej ustawia prompting 
jako element większego procesu, a nie "magiczny prompt", który wykona recenzję sam 
z siebie.
\subsection{Dostrajanie modeli (fine-tuning) i metody parameter-efficient}
\label{subsec:llm_acr_finetuning}

W podejściu opartym na dostrajaniu (\emph{fine-tuning}) przyjmuje się, że część ograniczeń promptowania 
(zwłaszcza to, że jakość bywa nierówna i czasami ciężko zapewnić oczekiwanej formy komentarza) da się ograniczyć
przez dopasowanie modelu do rozkładu danych pochodzących z procesu przeglądu kodu. Z punktu widzenia procesów
wytwarzania oprogramowania, wpasowuje się to w szerszą obserwację, że LLM w tym obszarze wykorzystuje się w wielu
zadaniach, lecz wykorzystanie bardzo ogólnego modelu do wyspecjalizowanego procesu, jakim jest code review,
zwykle wymaga dopasowania jednocześnie danych i celu uczenia \cite{Zhang2024SurveyLLMSE}.

\paragraph{Parameter-efficient fine-tuning jako kompromis wdrożeniowy.}
Dostrajanie oszczędne parametrowo (PEFT) jest jednym z ważniejszych kierunków w tym obszarze. Jest to podejście
w którym model adaptuje się do zadania bez modyfikowania pełnego zestawu wag. Lu i in. dobrze obrazują to podejście
poprzez LLaMA-Reviewer i pokazują, że zamiast trzymać osobne, ciężkie kopie modelu można przechowywać po prostu elementy dostrojenia.
Dla LLaMA 6.7B w ich opisie LoRA to ok. \(\sim 8.4\) mln trenowalnych parametrów i \(\sim 16\) MB, 
a prefix-tuning schodzi do ok. \(\sim 1.2\) mln parametrów i \(\sim 2.4\) MB \cite{Lu2023LLaMAReviewer}. PEFT nie jest
tam użyte tylko do jednego kroku, ale do kilku zadań w pipeline review (m.in. predykcji potrzeby review, 
generowania komentarza oraz \emph{code refinement}) \cite{Lu2023LLaMAReviewer}.
W tym miejscu istotne wydaje się takie wieloetapowie wykorzystanie tego podejścia, bo recenzja kodu w realnym narzędziu
zwykle nie kończy się na samym wygenerowaniu tekstu, tylko dochodzą kolejne etapy i decyzje. Podobny wniosek pojawia się
też w badaniach na zbiorze CodeReviewer: Haider i in. raportują, że PEFT (konkretnie QLoRA) daje mierzalne zyski w RCG, np. BLEU-4 rośnie do 5.58 
dla Code Llama 7B (vs 4.28 dla baseline CodeReviewer 223M), a BERTScore 
do 0.8483 dla Llama 3.1 8B (vs 0.8348) \cite{Haider2024PromptingFineTuningRCG}. Patrząc na to 
z perspektywy zdefiniowania state-of-the-art, wygląda na to, że fine-tuning nie jest "zarezerwowany" wyłącznie 
dla dużych laboratoriów i centr obliczeniowych, chociaż nadal łatwo tu o rozjazd wyników, jeśli dane 
albo protokół oceny są źle dobrane \cite{Tufano2024CodeReviewAutomation}.

\paragraph{Fine-tuning z naciskiem na zrozumiałość i strukturalną kompletność review.}
Faktem jest, że ocena jakości automatycznego code review nie powinna kończyć się na pytaniu "czy komentarz leksykalnie pasuje do referencji".
W praktyce w review liczy się raczej to, czy da się z komentarza wyciągnąć: gdzie jest problem, z czego on wynika i co można z 
tym zrobić (czyli jaka jest propozycja naprawy). Yu i in. projektują Carllm właśnie pod kątem poziomu jasności i
zrozumienia komentarzy i budują dane w oparciu o szerszy kontekst dyskusji review. Widać tu dość konkretne decyzje inżynierskie: z 1\,793 projektów 
open-source wybrano te spełniające kryteria aktywności (m.in. >1000 PR) oraz licencji (MIT/Apache 
2.0), a następnie z 15\,000 rekordów wyprowadzono 10\,554 (77.02\%) przykładów wysokiej jakości; 
po dalszym czyszczeniu zostało 9\,728 rekordów. Dalej dodano przykłady negatywne w proporcji 
1:1, co dało 19\,456 instancji treningowych \cite{Yu2024Carllm}. Autorzy zwracają też uwagę, że 
LoRA może obejmować bardzo mały ułamek parametrów (np. \(r=8\) to ok. 0.05\% 
parametrów modelu), więc sama adaptacja jest wykonalna nawet przy ograniczonych zasobach \cite{Yu2024Carllm}.
Wyniki benchmarkingowe pokazują istnienie ryzyka regresji (pogorszenia wynikowego kodu) 
nawet przy silnych modelach, więc nacisk na czytelność i sensowne uzasadnienie ma 
znaczenie nie tylko "użytkowe", ale też związane z bezpieczeństwem procesu \cite{Cihan2025EvaluatingLLMsForCodeReview}.

\paragraph{Jakość danych: destylacja "pożądanych" komentarzy jako warstwa kontrolna.}
Poprawa jakości samego sygnału uczącego jest dodatkowo ważnym wątkiem uzupełniającym. Yu i in. zwracają uwagę, że
komentarze review nie mają tej samej wartości, jeśli model ma się uczyć generowania uwag faktycznie prowadzących do
naprawy kodu. Autorzy wprowadzają pojęcie \emph{Desired Review Comments} (DRC) i metodę Desiview która przypisuje komentarzom
wynik pożądalności (\emph{desiredness score}). W praktyce jest to różnica perplexity (wariant z komentarzem vs bez komentarza), 
a wynik jest dodatkowo stabilizowany przez głosowanie czterech modeli (m.in. CodeLlama-13B, Meta-Llama-3-8B) 
\cite{Yu2025Desiview}. Na ręcznie oznaczanej próbce 600 rekordów Desiview osiąga 86.67\% accuracy oraz 
84.44\% F1 i wypada lepiej od GPT-4o (76.50\% accuracy, 71.05\% F1) \cite{Yu2025Desiview}. Autorzy pokazują też na przykładzie CodeReviewer,
że DRC to tylko ok. 43\% danych (np. 64\,934 z 150\,406 w treningu), więc "surowy" zbiór 
jest w dużej części zaszumiony, patrząc na cel generowania komentarzy prowadzących do 
realnych poprawek \cite{Yu2025Desiview}. Autorzy dostroili następnie model na wyczyszczonych danych, co poprawiło
BLEU-4 (np. dla LLaMA-3: 8.33 \(\rightarrow\) 11.87 po fine-tuningu, a po dodatkowym dopasowaniu do 13.13). Zwiększyło
też to odsetek ocen bardzo bliskich człowiekowi \cite{Yu2025Desiview}. Z perspektywy projektu wniosek jest prosty: samo dokładanie danych nie wystarczy,
a fine-tuning warto oprzeć na selekcji/filtracji przykładów, bo inaczej łatwo przepuścić do uczenia komentarze, które są mało 
pomocne.

\paragraph{Architektury wielozadaniowe i wykorzystanie MoE w procesie review.}
Tang i in. zwracają uwagę na problem: uczenie poszczególnych elementów pipeline'u 
recenzji kodu w izolacji sprawia, że wiedza nie przepływa między zadaniami tak, 
jak powinna. Chodzi głównie o relację między tym, czy system w ogóle uzna, że warto coś recenzować, a tym, jak potem faktycznie generuje 
dany komentarz. Wewnątrz mamy ekspertów wyspecjalizowanych w trzech konkretnych zadaniach: przewidywaniu konieczności recenzji (necessity 
prediction), generowaniu komentarzy oraz poprawianiu kodu (refinement). Całość spina mechanizm routowania Top-2\footnote{każdego wejścia wybierane są dwa najbardziej odpowiednie eksperty - modele, a ich wyniki są łączone}
wraz z load-balancingiem, co ma zapewniać odpowiednią efektywność \cite{Tang2025MoEReviewer}. 

Podejście to faktycznie ma dobre zastosowanie, co widać na wynikach na zbiorze CodeReviewer. Dla predykcji konieczności recenzji uzyskano 74,3\% accuracy (F1 na poziomie 
73,2\%). Przy samym generowaniu komentarzy BLEU wyniosło 11,62, a dla zadania refinement 
wynik ten dochodzi do 87,37 \cite{Tang2025MoEReviewer}. Bardzo ważnym aspektem tej pracy, zwłaszcza 
patrząc na to od strony czysto implementacyjnej czy bizneowej, jest wydajność obliczeniowa. Dzięki strukturze 
MoE model aktywuje tylko około 2 mld parametrów na token, co daje 
koszt rzędu 16,38 TFLOPs. To duży przeskok, bo wspomniany w pracy LLaMA-Reviewer 
(7B) wymaga aż 57,34 TFLOPs \cite{Tang2025MoEReviewer}. Ma to kluczowe znaczenie, bo narzędzia 
typu automated code review muszą przecież działać w ramach konkretnych budżetów czasowych i pieniężnych w procesach 
CI/CD. W takich warunkach opóźnienie inferencji i ogólny koszt są w zasadzie 
tak samo ważne, jak same metryki jakościowe modelu.

\paragraph{Wnioski z przeglądu literatury: fine-tuning a realia procesu review.}
W rzeczywistych warunkach code review jest procesem bardzo złożonym i wielowymiarowym. Sprawia to, że
ciężko jest sprowadzić to zadanie do prostego modelu 1:1, czyli relacji typu $diff \rightarrow$ pojedynczy komentarz.
Dobrze pokazują to Li i in. w pracy nad CodeDoctor - z ich analizy zbioru Turzo2023 wynika, że aż 78,7\% 
komentarzy dotyczyło co najmniej dwóch różnych kategorii. W przypadku zbioru Tufano2021 było 
to 56,71\%. Średnio na jeden komentarz przypadały dwie kategorie, a w ponad 
10\% przypadków było ich cztery lub nawet więcej \cite{Li2025CodeDoctor}. Ci sami autorzy 
zauważają, że w Turzo2023 zdecydowana większość procesów review (92,7\%) zajmowała programistom ponad 
pół godziny, co w zasadzie jest najlepszym argumentem za tym, żeby szukać 
metod automatyzacji wspierających recenzenta \cite{Li2025CodeDoctor}.

W pracach zbierających aktualny stan wiedzy, zauważyć można sporą dawkę sceptycyzmu wobec stosowania samych metryk.
Tufano i in. \cite{Tufano2024CodeReviewAutomation} trafnie wskazują, 
że bez dokładnej analizy jakościowej i kontroli danych łatwo wyciągnąć błędne wnioski. Bazowanie na samych metrykach
może ukrywać problemy z mapowaniem komentarzy na zmiany w kodzie. Może też maskować fakt, że dany model radzi sobie
bardzo dobrze tylko w tych najprostszych, powtarzalnych i często występujących przypadkach.

Bardzo ważnym aspektem jest także kwestia bezpieczeństwa oprogramowania. Germano i in. przeprowadzili systematyczny przegląd ponad dwustu publikacji z okresu 2018-sierpień 
2025 i widać tam ogromną dysproporcję w podejmowanych tematach. Aż 91,3\% prac 
skupia się wyłącznie na wykrywaniu podatności Z kolei naprawianie błędów
czy ich wyjaśnianie to odpowiednio tylko 11,1\% i 5,3\% publikacji \cite{Germano2025VulnSLR}. 
To pokazuje lukę w badaniach - w systemach automatycznego przeglądu kodu nie wystarczy przecież 
samo "wskazanie palcem" problemu. Żeby narzędzie było użyteczne, musi oferować zrozumiałe uzasadnienie 
i konkretne propozycje zmian, co zresztą potwierdzają autorzy prac nad podejściami DRC 
i comprehensibility \cite{Yu2024Carllm,Yu2025Desiview}.

Wdrożenie w przemyśle jest równie istotnym aspektem. Ciekawym przykładem są doświadczenia z firmy Ericsson, opisane przez 
Ramesha \cite{Ramesh2025EricssonExperience}. Tam zamiast kosztownego fine-tuningu postawiono na "lekkie" podejście zintegrowane bezpośrednio 
z workflow (Gerrit, VS Code). Wykorzystali oni parser Tree-sitter do wyciągania kontekstu 
z AST (tzw. enclosing method), co pokazuje, że w praktyce wybór modelu to zawsze kompromis. 
Wynika z tego, że w praktyce należy wyważyć jakość wyników, koszty utrzymania 
i to, jak trudno będzie zintegrować narzędzie z istniejącym procesem wytwórczym \cite{Ramesh2025EricssonExperience}.

\paragraph{Podsumowanie.}
Dostrajanie modeli (szczególnie PEFT) jest dziś jednym z ważniejszych elementów state-of-the-art w 
automatycznym \emph{code review}. Pozwala utrzymać bardziej przewidywalny styl i dość 
spójną strukturę komentarzy. Pozwala też lepiej dopasować zachowanie modelu do tego, jak 
wyglądają dane w konkretnym repozytorium. Ekuteczność takiego podejścia mocno zależy od: jego jakości, oraz udziale komentarzy, które 
faktycznie uznajemy za pożądane. Do tego dochodzi kwestia tego, jak w ogóle 
mierzymy jakość (protokół ewaluacji) oraz koszty wdrożenia, które czasem są niesłusznie pomijane.

Część typowych problemów w \emph{review} (braki kontekstu projektowego, halucynacje albo zgodność z regułami jakości) bywa
łatwiej ograniczać metodami hybrydowymi i/lub przez retrieval, zamiast samym dostrajaniem. W literaturze pojawia się to m.in. jako klasyczny 
RAG dopasowany do \emph{review} \cite{Meng2025RARe}, ale też jako integracja z analizą statyczną 
(RAG oparty o wyniki analizatorów) \cite{Jaoua2025LLMStaticAnalyzersCR} albo jako dostarczanie kuratorowanej "mapy wiedzy", 
która działa bardziej jak element rozumowania symbolicznego \cite{Icoz2025SymbolicReasoning}. W skrócie: PEFT pomaga, 
tylko nie rozwiązuje wszystkiego sam, a w części przypadków sensownie jest go 
łączyć z retrieval lub dodatkowymi źródłami sygnału.

\subsection{Kontekst w code review i systemy hybrydowe: RAG, statyka i logika}
\label{subsec:llm_acr_context_hybrid}

Modele LLM posiadają jeden główny problem w kontekście automatycznego code review, niezależnie czy
korzystamy z gotowych modeli czy dodajemy douczenie - modele wiedzą za mało. W realnym procesie przeglądu kodu
programista bardzo rzadko opiera się wyłącznie na zaznaczonych zmianach w kodzie. Równie ważne, lub nawet ważniejsze jest
posiadanie wiedzy na temat konwencji przyjętych w projekcie, architektury, historii poprzednich zmian czy zależności
pomiędzy plikami, metodami czy modułami. Istotne są też dane z testów, statycznej analizy czy linterów. Są to rzeczy
których brakuje w podstawowym podejściu użycia modelu językowego. Z tego powodu powstał trend polegający na nietraktowaniu
LLM jako jedynego i samodzielnego sędziego, który orzeka o poprawności kodu. Zamiast tego próbuje się go osadzić w szerszych
systemach, gdzie dostarczana jest mu wiedza zewnętrzna i inne narzędzia pilnują jego wyników. O tych problemach
pisze Tufano i inni w swoim przeglądzie \cite{Tufano2024CodeReviewAutomation}. Zwracają oni uwagę na to, że automatyzacja review ma sporo słabych punktów, jak 
choćby to, że metryki tekstowe bywają mało wiarygodne, a samo mapowanie komentarzy do konkretnych 
zmian w kodzie bywa po prostu niejasne, co utrudnia wyciąganie sensownych wniosków.

\paragraph{Zastosowanie Retrieval-Augmented Generation jako sposobu na dostarczenie kontekstu projektowego}
Mechanizm \emph{retrieval-augmented generation} (RAG) jest podejściem, które pozwala na poradzenie sobie z problemem
braku odpowiedniego kontekstu. W tym modelu przed generowaniem komentarzy, przeszukuje repozytorium w poszukiwaniu
powiązanych elementów: wcześniejsze dyskusje, dokumentacja, podobne zmiany w kodzie. Tak przygotowane dane
trafiają do modelu jako ustrukturyzowany kontekst. Meng i in. \cite{Meng2025RARe} trafnie wskazują tutaj na kluczowy aspekt:
samą automatyzację review warto opierać na bardzo precyzyjnej selekcji i "dawkowaniu" tych informacji. Celem jest to,
żeby zwiększyć trafność podpowiedzi bez przeładowania promptu niepotrzebnymi danymi. Finalna jakość zależy właściwie w
równym stopniu od samego modelu LLM, co od sprawności wyszukiwania, czyli tego, jak dobierzemy źródła, reprezentacje danych czy same algorytmy szeregowania 
i deduplikacji tych fragmentów.

RAG w kontekście tworzonego projektu można potraktować jako swoisty most, który łączy konkretne zmiany w 
PR/MR z potrzebną wiedzą zgromadzoną w repozytorium. Daje to szansę na to, by komentarze odnosiły się do faktycznych
standardów panujących w projekcie (widząc np. jak podobne problemy rozwiązano wcześniej). Sprawia to też, że możliwe jest
ograniczenie generowanie pustych ogólników, gdzie jest to częsty problem wytykany automatycznym metodom code review \cite{Ren2025HydraReviewer}.
Jest to bardzo ważne w systemach, które mają realnie wspierać recenzenta — system powinien podsuwać mu konkretne, łatwe do 
zweryfikowania argumenty, zamiast seryjnie produkować generyczne uwagi.

\paragraph{Integracja z analizą statyczną: hybryda precyzji i elastyczności.}
Łączenie LLM z analizatorami statycznymi jest drugim nurtem w podejściach hybrydowych. Chodzi w tym podejściu o to,
że klasyczne statyczne analizatory kodu są zwykle bardzo precyzyjne w znajdywaniu konkretnych klas problemów w kodzie,
ale ich raporty są mało przyjazne i nie określają sposobu czy propozycji naprawy. Często brakuje w nich kontekstu, co dokładnie jest nie 
tak z punktu widzenia projektu i dlaczego to ma znaczenie. LLM może wtedy przełożyć sygnał narzędziowy na czytelny komentarz.
Ważne jest to, że inicjacja nadal opiera się na detekcjach zewnętrznych. Taki kierunek opisują autorzy pracy 
o łączeniu LLM z analizatorami statycznymi do generowania review, gdzie wyniki analiz pełnią rolę 
dodatkowego kontekstu i ograniczeń dla generacji \cite{Jaoua2025LLMStaticAnalyzersCR}. Od strony kontroli jakości to jest zachęcające podejście -
da się wprost pokazać, że komentarz wynika z konkretnej reguły albo ostrzeżenia narzędzia.

Pasuje to do obserwacji dotyczących ryzyka regresji i ogólnej niepewności rekomendacji modelu. Skoro nawet zaawansowane LLM
potrafią zaproponować zmianę, która pogarsza poprawny kod, to umocnienie procesu w wynikach analiz oraz walidacji w pipeline może działać
jak bufor bezpieczeństwa \cite{Cihan2025EvaluatingLLMsForCodeReview}. Zatem analiza statyczna nie jest więc tylko "źródłem wiedzy",
ale też elementem, co zawęża przestrzeń dopuszczalnych sugestii i wymusza bardziej sprawdzalne uzasadnienie.

\paragraph{Podejścia hybrydowe i wykorzystanie logiki symbolicznej.} 
Istotnym kierunkiem jest włączanie w proces generowania treści elementów symbolicznych. Przykładowo, w pracy
\cite{Icoz2025SymbolicReasoning} autorzy opisują system do automatycznego review, w którym LLM współpracuje z warstwą regułową opartą o tzw. mapę wiedzy.
Takie połączenie ma za zadanie pilnować spójności i ograniczać sytuacje, w których model generuje błędy, które na pierwszy 
rzut oka brzmią bardzo sensownie (tzw. halucynacje). Z punktu widzenia programisty takie podejście jest 
ciekawe głównie ze względu na faktyczną możliwość wyjaśnienia - dużo łatwiej zaakceptować komentarz do kodu, jeśli 
system potrafi go uzasadnić konkretną regułą, którą da się zweryfikować, a nie opiera się 
tylko na "intuicji" modelu.

Łączenie logiki symbolicznej z modelami językowymi ma spore znaczenie, jeśli chodzi o bezpieczeństwo i analizę podatności w kodzie.
Systematyczny przegląd literatury przeprowadzony w \cite{Germano2025VulnSLR} pokazuje jednak 
pewną lukę w obecnych badaniach. Sama detekcja podatności przy użyciu LLM jest już dość szeroko opisana, ale kwestie automatycznej
naprawy czy wyjaśniania, skąd dany problem się wziął, są traktowane jako rzeczy o niskim priorytecie. Podczas tworzenia
narzędzia do automatycznego code review jest to ważny aspekt - samo oznaczenie fragmentu kodu jako problematyczny to często za mało.
Wychodzi z tego, że to właśnie mechanizmy hybrydowe które potrafią odpowiedzieć na pytania "dlaczego" i "jak to naprawić", będą 
kluczowe, żeby generowane komentarze faktycznie poprawiały jakość oprogramowania, a nie tylko zwiększały liczbę powiadomień 
dla dewelopera.

\paragraph{Wnioski z praktyki przemysłowej: kontekst strukturalny i integracja workflow.}
Dobry model sam z siebie nie da sensownych efektów, jeśli nie jest sensownie zintegrowany w proces wytwarzania oprogramowania.
Pokazują to przemysłowe badania, w tym doświadczenia z Ericsson, gdzie pokazano, że w praktyce sporo poświęcone jest na integrację z narzędziami, z których i 
tak korzystają programiści (np. środowisko review oraz edytor), a także na wyciąganie kontekstu strukturalnego 
kodu (np. metody otaczającej zmianę). To drugie jest o tyle istotne, że zmniejsza liczbę sytuacji, w których model ocenia fragment kodu jakby był wyizolowany,
bez związku z jego rolą w programie \cite{Ramesh2025EricssonExperience}. Jednak podejście zwiększające ilość kontekstu
powinno być przemyślane. Wygrywa tutaj lepszy, bardziej semantyczny opis kontekstu 
(np. AST, zależności, lokalizacja zmiany), nawet jeśli czasem wygląda to mniej efektownie na wejściu.

\paragraph{Podsumowanie.}
Wzbogacanie kontekstu oraz podejścia hybrydowe stają się jednym z ważniejszych kierunków rozwoju automatycznego \emph{code review},
jak pokazuje podany zbiór literatury. RAG pozwala na dołożenie do modelu informacji projektowych i dzięki temu możliwe jest
ograniczenie zbyt ogólnych lub nietrafnych komentarzy \cite{Meng2025RARe,Ren2025HydraReviewer}.
Integracja z analizą statyczną zwykle poprawia weryfikowalność (łatwiej sprawdzić, skąd bierze się sugestia) i może również zmniejszać ryzyko nietrafionych 
rekomendacji \cite{Jaoua2025LLMStaticAnalyzersCR,Cihan2025EvaluatingLLMsForCodeReview}.
Dołożenie komponentów symbolicznych pomaga utrzymać spójność rozumowania i daje lepszą wyjaśnialność, chociaż nie zawsze jest to tanie pod kątem 
złożoności systemu \cite{Icoz2025SymbolicReasoning}. Coraz częściej spotyka się narzędzia do automatycznego code review jako kompozycję: LLM działa 
jako generator, a zewnętrzne źródła wiedzy oraz narzędzia pełnią rolę ugruntowania i kontroli jakości \cite{Tufano2024CodeReviewAutomation,Ramesh2025EricssonExperience}.

\subsection{Ewaluacja, ograniczenia i luki badawcze w automatycznym code review z użyciem LLM}
\label{subsec:llm_acr_evaluation_gaps}

Ocena tego, jak systemy ACR oparte na modelach językowych radzą sobie w praktyce, 
to w zasadzie jeden z najtrudniejszych tematów w literaturze. Wynika to z faktu, że code review to też proces silnie
społeczny. Tak jak już zauważono w sekcji \ref{subsec:llm_acr_prompting}, poleganie wyłącznie na miarach
ilościowych bywa mało miarodajne. Tufano i inni \cite{Tufano2024CodeReviewAutomation} trafnie punktują tutaj kwestię przygotowania zbiorów 
danych - jeśli mapowanie komentarzy do konkretnych zmian w kodzie jest nieprecyzyjne, to wyniki 
modelu mogą wyglądać na papierze bardzo dobrze, podczas gdy w realnym procesie pracy programisty będą 
zupełnie nieprzydatne. Prowadzi to do co raz częściej występujących głosów w literaturze, że konieczne jest używanie
również weryfikacji funkcjonalnej obok metryk automatycznych \cite{Tufano2024CodeReviewAutomation, Zhang2024SurveyLLMSE}.
Nie jest to jednak proste do wdrożenia w badaniach.

\paragraph{Metryki tekstowe a poprawność merytoryczna.}
Metryki generacji tekstu, takie jak BLEU czy BERTScore to kolejny aspekt warty omówienia. Haider 
i in. \cite{Haider2024PromptingFineTuningRCG} pokazali, że pozwalają one na pewną formę porównywania wyników między modelami, 
to tak naprawdę niewiele mówią one o tym, czy wygenerowana porada ma sens merytoryczny i czy jest zgodna z konkretnymi
standardami danego projektu \cite{Tufano2024CodeReviewAutomation}. Problem nie jest łatwy w rozwiązaniu, ponieważ
komentarze w procesie przeglądu kodu dotyczą skrajnie różnych rzeczy: czystości stylu, wydajność kodu, poprawność
logiki biznesowej i inne. Ciężko więc oczekiwać, że jakość recenzji da się 
sprowadzić do jakiegoś jednego, uniwersalnego wskaźnika liczbowego, co zresztą potwierdzają też autorzy pracy \cite{Li2025CodeDoctor}, 
wskazując na trudności w obiektywnej kategoryzacji takich uwag.

\paragraph{Ewaluacja funkcjonalna i ryzyko regresji.}
Praktyka bardziej skupiona na odpowiedzi, czy sugerowana zmiana faktycznie poprawia
kod, a nie tylko wygląda "rozsądnie" na papierze jest bliższe rozsądkowi. Protokół walidacji oparty o testy jednostkowe, zaproponowany
przez Cihana i in. i omówiony wcześniej przy okazji intencji modelu, pokazuje przy tym dość jasno, że ryzyko wprowadzenia
zmian szkodliwych nie jest zerowe \cite{Cihan2025EvaluatingLLMsForCodeReview}. W kontekście projektowanego systemu
oznacza to konieczność dołożenia elementów zabezpieczających, np. automatycznej walidacji w procesach CI
oraz podejścia \emph{human-in-the-loop}.

\paragraph{Ewaluacja diagnostyczna i znaczenie kontekstu.}
Faktycznie zrozumienie zmian w code review jest nowoczesnym podejściem które co raz częściej widać w systemach automatyzacji
tego procesu. Służą do tego dedykowane benchmarki, takie jak CodeReviewQA 
\cite{Lin2025CodeReviewQA}, które pozwalają sprawdzić, czy model w ogóle potrafi poprawnie zinterpretować intencje stojące za 
zmianami w kodzie, zanim jeszcze zaproponuje konkretne poprawki \cite{Lin2025CodeReviewQA}. Odpowiednie osadzenie modelu
w kontekście konkretnego repozytorium do kolejny istotny aspekt widoczny w literaturze. Zastosowanie mechanizmów RAG albo
łączenie LLM-ów z analizą statyczną — co opisują Meng \cite{Meng2025RARe} oraz Jaoua \cite{Jaoua2025LLMStaticAnalyzersCR} — nie tylko
poprawia samą jakość generowanych odpowiedzi, ale też sprawia, że łatwiej jest prześledzić ich pochodzenie.

\paragraph{Luki w badaniach i wnioski dla projektowania systemów.}
Przy analizowaniu obecnych rozwiązań ACR widać jednak kilka wyraźnych braków:
\begin{itemize}
    \item \textbf{Niepełny cykl pracy}: Większość prac skupia się właściwie tylko na tym, 
jak wygenerować komentarz do kodu. Jak zauważa Germano \cite{Germano2025VulnSLR}, brakuje tutaj szerszego spojrzenia, które 
obejmowałoby też etap uzasadniania decyzji czy późniejszą, automatyczną naprawę wykrytego problemu.
    \item \textbf{Problemy z jakością danych treningowych}: Same dane bywają problematyczne. Yu i 
in. \cite{Yu2025Desiview} słusznie punktują, że bez rygorystycznego filtrowania (np. poprzez destylację DRC), modele po 
prostu uczą się od ludzi mało efektywnych wzorców komunikacji, co potem przekłada się na 
słabą jakość sugestii.
    \item \textbf{Praktyczne bariery przy wdrożeniach}: Na koniec zostają kwestie czysto techniczne i 
biznesowe. Doświadczenia opisane przez Ramesha na przykładzie Ericssona \cite{Ramesh2025EricssonExperience} pokazują, że finalnie to nie 
parametry modelu, a koszty, czas odpowiedzi i trudność zintegrowania narzędzia z obecnym workflow programistów 
decydują o tym, czy system się przyjmie.
\end{itemize}

Biorąc powyższe pod uwagę, stwierdzić można, że budowa efektywnego narzędzia do wspomagania przeglądu kodu
wymaga więcej niż prostego generowania kodu. Kluczowe okazuje się tutaj 
połączenie selektywnego podejścia do kontekstu z bardzo dokładną, wielowarstwową 
procedurą ewaluacji. Ocena nie może ograniczać się tylko do samej poprawności składni; 
musi brać pod uwagę też realną użyteczność i pilnować, żeby nie dochodziło 
do regresji w kodzie \cite{Tufano2024CodeReviewAutomation, Zhang2024SurveyLLMSE}. 
Wspomniane aspekty są kluczowe dla zapewnienia stabilności narzędzia.
\subsection{Sposoby oceny jakości automatycznego code review}
\label{subsec:llm_acr_eval}

Właściwe mierzenie skuteczności systemów do automatycznego przeglądu kodu jest problematyczną kwestią.
Największym problemem jest to, że klasyczne miary używane do przetwarzania języka naturalnego rzadko kiedy
oddają to, czy komentarz code review jest faktycznie przydatny technicznie \cite{Lin2025CodeReviewQA, Zhang2024SurveyLLMSE}.
W publikacjach widać trzy główne podejścia do ocen: metryki ilościowe, badania jakościowe z użyciem programistów oraz
gotowe benchmarki badawcze.

\paragraph{Metryki ilościowe i ich ograniczenia}
W zadaniach typu \textit{code-to-comment} czy przy automatycznych poprawkach (\textit{code refinement}), badacze zazwyczaj sięgają po 
metryki leksykalne - mowa tu głównie o BLEU-4, ROUGE-L oraz METEOR \cite{Lin2025CodeReviewQA, Meng2025RARe, Yu2024Carllm}. 
Mimo ich powszechnego użycia, wiarygodność tych miar jest coraz częściej podważana w nowszych pracach:

\begin{itemize}
    \item \textbf{Krytyka Exact Match (EM):} Ta metryka bywa oceniana jako zbyt restrykcyjna, 
zwłaszcza w zadaniach, gdzie liczy się sens, a nie identyczne brzmienie. 
W pracy Tufano \cite{Tufano2024CodeReviewAutomation} ręczna analiza pokazała, że nawet 21,83\% predykcji odrzuconych przez 
EM w zadaniu generowania komentarzy było jednak semantycznie poprawnych. To dość mocno sugeruje, że samo EM może 
"karać" poprawne odpowiedzi tylko dlatego, że są sformułowane inaczej. 
W pracy Pornprasit \cite{Pornprasit2024FineTuningPromptingCR} pojawia się podobny wniosek, gdzie zwraca się uwagę, że 
EM może zaniżać realne możliwości modeli w porównaniu z oceną ludzką (która jest bardziej elastyczna).
    
    \item \textbf{BERTScore:} Jest to podejście oparte o aspekt kontekstowy, więc zamiast patrzeć 
na identyczność tokenów, próbuje uchwycić podobieństwo znaczenia. Dla komentarzy do code review ma to sens, bo pozwala lepiej 
oddać intencję recenzenta niż proste dopasowanie słów, co podkreślają autorzy pracy Haider \cite{Haider2024PromptingFineTuningRCG}.
    
    \item \textbf{Desiredness Score (DS):} Opiera się na różnicy \textit{perplexity} modelu.
Zamysł jest taki, że mierzy się, na ile dany komentarz zmniejsza niepewność modelu przy generowaniu poprawnej wersji kodu.
Dzięki temu da się automatycznie wyłapywać tzw. \textit{Desired Review Comments} (DRC), czyli takie uwagi, które faktycznie prowadzą do 
napraw, a nie tylko "brzmią sensownie" \cite{Yu2025Desiview}.
    
    \item \textbf{Wskaźniki poprawności technicznej:} Coraz częściej podkreśla się potrzebę mierzenia nie tylko 
jakości językowej komentarzy, ale też efektu w kodzie. W szczególności istotne są \textit{Correction Ratio} (czy poprawki realnie działają)
oraz \textit{Regression Ratio}, które opisuje ryzyko, że sugestia modelu popsuje wcześniej poprawny
kod \cite{Cihan2025EvaluatingLLMsForCodeReview}. To jest o tyle ważne, że nawet "ładna" uwaga w code review nie ma wartości,
jeśli w praktyce prowadzi do regresji.
\end{itemize}

\paragraph{Oceny eksperckie i badania użyteczności}
Podczas typowego wdrożenia przemysłowego automatyczne metrki to z reguły za mało, więc dalej kluczową rolę w procesach
code review odgrywa przegląd przez doświadczonego programistę \cite{Ramesh2025EricssonExperience}. 
W nowych badaniach dodatkowo zaczęto odchodzić od ogólnych ocen na rzecz wielowymiarowych skal.
W takim podejściu zwraca się uwagę na trafność (\textit{relevance}), czyli to, czy dany komentarz w 
ogóle odnosi się do właściwego fragmentu kodu i czy poprawnie wskazuje źródło problemu \cite{Yu2024Carllm, Ren2025HydraReviewer}.
Kolejnym jest użyteczność (\textit{actionability}) - czy sugestia jest na tyle konkretna, żeby programista wiedział, jak ma wprowadzić poprawkę.
Jest to o tyle istotne, że według badań Ren \cite{Ren2025HydraReviewer} aż 61,0\% uwag generowanych przez tradycyjne metody 
jest ocenianych jako zbyt ogólnikowe (\textit{vagueness}), co czyni je bezużytecznymi.
Ponad to jest kwestia zrozumiałości (\textit{comprehensibility}) - czy wyjaśnienie błędu jest jasne dla człowieka i
czy model, generując odpowiedź, zachował logiczny ciąg rozumowania, co często wiąże się z zastosowaniem
techniki \textit{Chain-of-Thought} \cite{Yu2024Carllm, Haider2024PromptingFineTuningRCG}.

Istnieje także podejście \textit{LLM-as-a-judge}. Jest to wykorzystanie zaawansowanych i większych modeli do oceny
wyników pracy mniejszych systemów. Wyniki takich eksperymentów wyglądają obiecująco — w pracy Jaoua \cite{Jaoua2025LLMStaticAnalyzersCR} wykazano, 
że taka automatyczna ocena wykazuje dość wysoką zgodność z opiniami ludzi, osiągając współczynnik \textit{Cohen's 
kappa} na poziomie 0,72.

\paragraph{Benchmarki i dekompozycja zadań}
Przegląd kodu co raz rzadziej jest traktowany jako jedno, monolityczne zadanie. Częściej 
rozbija się je na etapy odpowiadające kolejnym krokom rozumowania. Przykładem jest benchmark \textit{CodeReviewQA}, w 
którym sprawdza się działanie modelu w trzech krokach \cite{Lin2025CodeReviewQA}:
\begin{enumerate}
    \item \textbf{CTR (\textit{Change Type Recognition}):} rozpoznanie, jaki typ zmiany jest potrzebny.
    \item \textbf{CL (\textit{Change Localisation}):} wskazanie konkretnych linii kodu, których dotyczy problem. W 
praktyce to właśnie lokalizacja bywa najczęstszym "wąskim gardłem" dla modeli LLM \cite{Lin2025CodeReviewQA}.
    \item \textbf{SI (\textit{Solution Identification}):} wybór właściwej poprawki spośród dostępnych alternatyw.
\end{enumerate}

Trzeba dodatkowo brać pod uwagę problemy metodologiczne. Dane historyczne 
potrafią być zaszumione: w popularnych zbiorach szum szacuje się na około 25\% rekordów \cite{Tufano2024CodeReviewAutomation}.
Istnieje ryzyko \textit{data leakage} (wycieku danych), które w zadaniach naprawy kodu 
potrafi zawyżać wyniki o ponad 30 punktów procentowych \cite{Germano2025VulnSLR}.


